% ---------------------------------------------------------------------
\subsubsection{Requisiti Funzionali}
\label{subsec:requisiti-funzionali}
% ---------------------------------------------------------------------

Questa sottosezione descrive in dettaglio tutte le funzionalità del sistema \productname, specificando input, elaborazione e output secondo lo standard IEEE 830-1998. Il sistema è organizzato gerarchicamente: gli \textbf{Amministratori} creano \textbf{Progetti} e assegnano \textbf{Utenti} ad essi; all'interno di ogni progetto, i membri autorizzati gestiscono \textbf{Issue}.

\vspace{0.5cm}

% =====================================================================
% REQ-FUN-01
% =====================================================================

\begin{reqbox}[title=REQ-FUN-01: Autenticazione]
\label{subsubsec:req-autenticazione}

\subparagraph{Introduzione}
Il sistema deve fornire un meccanismo di autenticazione sicuro per accesso alla piattaforma.

\subparagraph{Input}
\begin{itemize}[leftmargin=1.1em]
\item Email (obbligatorio, formato valido)
\item Password (obbligatorio, min 8 caratteri)
\end{itemize}

\subparagraph{Elaborazione}
\begin{enumerate}[leftmargin=1.1em]
\item Valida formato email e presenza campi
\item Cerca account nel database
\item Se valida: crea sessione
\item Se non valida: incrementa contatore tentativi
\end{enumerate}

\subparagraph{Output}
\textbf{Successo}: Reindirizzamento a dashboard progetti

\textbf{Errore}: "Credenziali non valide", "Account bloccato", "Campi obbligatori"

\end{reqbox}

\vspace{0.5cm}

\newpage

% =====================================================================
% REQ-FUN-02
% =====================================================================

\begin{reqbox}[title=REQ-FUN-02: Visualizza Dashboard Progetti]
\label{subsubsec:req-dashboard-progetti}

\subparagraph{Introduzione}
Dopo l'autenticazione, l'utente visualizza la dashboard con i \textbf{progetti a cui ha accesso} assegnati dall'amministratore.

\subparagraph{Precondizioni}
Utente autenticato.

\subparagraph{Input}
Caricamento automatico post-login.

\subparagraph{Elaborazione}
\begin{enumerate}[leftmargin=1.1em]
\item Recupera l'utente autenticato dal database
\item Recupera i progetti legati all'utente
\end{enumerate}

\subparagraph{Output}
Dashboard con:
\begin{itemize}[leftmargin=1.1em]
\item Card progetti: nome, descrizione, statistiche
\item Pulsante "Accedi" per selezionare un progetto
\item Se nessun progetto: "Nessun progetto disponibile."
\item Se admin: pulsante "+ Crea Nuovo Progetto"
\end{itemize}

\end{reqbox}

\vspace{0.5cm}

\newpage

% =====================================================================
% REQ-FUN-03
% =====================================================================

\begin{reqbox}[title=REQ-FUN-03: Seleziona Progetto]
\label{subsubsec:req-seleziona-progetto}

\paragraph{Introduzione}
L'utente seleziona un progetto per accedere alle sue issue.

\paragraph{Precondizioni}
Utente autenticato con almeno un progetto assegnato.

\paragraph{Input}
Click su progetto dalla dashboard.

\paragraph{Elaborazione}
\begin{enumerate}[leftmargin=1.1em]
\item Verifica che l'utente appartenga al progetto selezionato
\item Imposta il progetto corrente nella sessione
\item Recupera le issue del progetto
\item Reindirizza alla lista issue
\end{enumerate}

\paragraph{Output}
\textbf{Successo}: Entra nel contesto del progetto, header mostra nome progetto, visualizza lista issue

\textbf{Errore}: "Accesso negato. Non fai parte di questo progetto."

\end{reqbox}

\vspace{0.5cm}

\newpage

% =====================================================================
% REQ-FUN-04
% =====================================================================

\begin{reqbox}[title=REQ-FUN-04: Crea Progetto (solo Admin)]
\label{subsubsec:req-crea-progetto}

\paragraph{Introduzione}
Solo gli amministratori possono creare nuovi progetti.

\paragraph{Precondizioni}
Utente con ruolo Amministratore.

\paragraph{Input}
\begin{itemize}[leftmargin=1.1em]
\item Nome progetto (obbligatorio, univoco)
\item Descrizione (opzionale)
\item Membri iniziali (obbligatorio, multi-selezione utenti, con almeno un membro)
\end{itemize}

\paragraph{Elaborazione}
\begin{enumerate}[leftmargin=1.1em]
\item Verifica ruolo amministratore
\item Valida unicità nome progetto
\item Crea nuovo progetto nel database
\item Associa i membri selezionati al progetto
\item Aggiunge automaticamente l'amministratore creatore come membro
\end{enumerate}

\paragraph{Output}
\textbf{Successo}: "Progetto '[Nome]' creato con successo"

\textbf{Errore}: "Nome già esistente", "Permessi insufficienti"

\end{reqbox}

\vspace{0.5cm}

\newpage

% =====================================================================
% REQ-FUN-05
% =====================================================================

\begin{reqbox}[title=REQ-FUN-05: Modifica Progetto (solo Admin)]
\label{subsubsec:req-modifica-progetto}

\paragraph{Introduzione}
Gli amministratori possono modificare nome e descrizione di progetti esistenti.

\paragraph{Precondizioni}
Amministratore, progetto esistente.

\paragraph{Input}
Nuovo nome (se modificato deve essere univoco), nuova descrizione.

\paragraph{Elaborazione}
\begin{enumerate}[leftmargin=1.1em]
\item Verifica ruolo amministratore
\item Se nome modificato: verifica unicità
\item Aggiorna i dati del progetto nel database
\end{enumerate}

\paragraph{Output}
\textbf{Successo}: "Progetto aggiornato con successo"

\textbf{Errore}: "Nome già in uso", "Permessi insufficienti"

\end{reqbox}

\vspace{0.5cm}

\newpage

% =====================================================================
% REQ-FUN-06
% =====================================================================

\begin{reqbox}[title=REQ-FUN-06: Elimina Progetto (solo Admin)]
\label{subsubsec:req-elimina-progetto}

\paragraph{Introduzione}
Gli amministratori possono eliminare progetti con tutte le issue associate.

\paragraph{Precondizioni}
Amministratore, progetto esistente.

\paragraph{Input}
Selezione progetto, conferma con password amministratore.

\paragraph{Elaborazione}
\begin{enumerate}[leftmargin=1.1em]
\item Mostra dialog conferma: "Eliminare progetto? Tutte le issue saranno eliminate. Operazione irreversibile."
\item Verifica password amministratore
\item Elimina tutte le issue, commenti, etichette e allegati del progetto
\item Rimuove associazioni membri-progetto
\item Elimina il progetto dal database
\end{enumerate}

\paragraph{Output}
\textbf{Successo}: "Progetto '[Nome]' e tutte le sue issue eliminate"

\textbf{Errore}: "Password non valida"

\end{reqbox}

\vspace{0.5cm}

\newpage

% =====================================================================
% REQ-FUN-07
% =====================================================================

\begin{reqbox}[title=REQ-FUN-07: Aggiungi Membri Progetto (solo Admin)]
\label{subsubsec:req-gestisci-membri}

\paragraph{Introduzione}
Gli amministratori decidono quali utenti possono accedere a ciascun progetto.

\paragraph{Precondizioni}
Amministratore, progetto esistente.

\paragraph{Input}
Selezione utenti da aggiungere al progetto.

\paragraph{Elaborazione}

\textbf{Aggiunta membro}:
\begin{enumerate}[leftmargin=1.1em]
\item Verifica che l'utente non sia già membro del progetto
\item Associa l'utente al progetto nel database
\item Notifica l'utente dell'aggiunta
\end{enumerate}

\paragraph{Output}
\textbf{Successo}: "Utente [Nome] aggiunto al progetto"

\textbf{Errore}: "Utente già membro".

\end{reqbox}

\vspace{0.5cm}

\newpage

% =====================================================================
% REQ-FUN-08
% =====================================================================

\begin{reqbox}[title=REQ-FUN-08: Rimuovi Membri Progetto (solo Admin)]
\label{subsubsec:req-gestisci-membri}

\paragraph{Introduzione}
Gli amministratori decidono quali utenti devono essere rimossi da un determinato progetto.

\paragraph{Precondizioni}
Amministratore, progetto esistente, utente\_autenticato membro del progetto

\paragraph{Input}
Selezione utenti da rimuovere dal progetto.

\paragraph{Elaborazione}

\textbf{Rimozione membro}:
\begin{enumerate}[leftmargin=1.1em]
\item Verifica che il progetto abbia almeno un altro membro
\item Rimuove le assegnazioni issue dell'utente nel progetto
\item Rimuove l'associazione utente-progetto dal database
\item Notifica l'utente della rimozione
\end{enumerate}

\paragraph{Output}
\textbf{Successo}: "Utente [Nome] rimosso dal progetto"

\textbf{Errore}: "Impossibile rimuovere ultimo membro"

\end{reqbox}

\vspace{0.5cm}

\newpage

% =====================================================================
% REQ-FUN-09
% =====================================================================

\begin{reqbox}[title=REQ-FUN-09: Segnala Issue]
\label{subsubsec:req-segnala-issue}

\paragraph{Introduzione}
Gli utenti creano nuove issue all'interno del \textbf{progetto correntemente selezionato}.

\paragraph{Precondizioni}
Utente autenticato, progetto selezionato, utente membro del progetto.

\paragraph{Input}
\begin{itemize}[leftmargin=1.1em]
\item Titolo (obbligatorio)
\item Descrizione (obbligatoria)
\item Tipo (obbligatorio: Bug, Feature, Question, Documentation)
\item Priorità (opzionale: Nessuna, Bassa, Media, Alta, Critica)
\item Etichette (opzionale)
\item Immagini (opzionale)
\end{itemize}

\paragraph{Elaborazione}
\begin{enumerate}[leftmargin=1.1em]
\item Verifica che l'utente appartenga al progetto corrente
\item Valida tutti i campi obbligatori e i vincoli di formato
\item Se presenti immagini: carica i file su storage
\item Se priorità non specificata: imposta a "Nessuna"
\item Crea la nuova issue associata al progetto corrente con stato "ToDo"
\item Se priorità è "Critica": notifica gli amministratori del progetto
\end{enumerate}

\paragraph{Output}
\textbf{Successo}: "Issue \#[ID] creata con successo", reindirizzamento a dettaglio issue

\textbf{Errore}: "Campo obbligatorio mancante", "Immagine troppo grande", "Formato non supportato"

\end{reqbox}

\vspace{0.5cm}

\newpage

% =====================================================================
% REQ-FUN-10
% =====================================================================

\begin{reqbox}[title=REQ-FUN-10: Visualizza Lista Issue]
\label{subsubsec:req-visualizza-lista-issue}

\paragraph{Introduzione}
Gli utenti visualizzano tutte le issue del \textbf{progetto corrente}.

\paragraph{Precondizioni}
Utente autenticato, progetto selezionato, membro del progetto.

\paragraph{Input}
Accesso a sezione "Issue" o automatico dopo selezione progetto.

\paragraph{Elaborazione}
\begin{enumerate}[leftmargin=1.1em]
\item Verifica appartenenza al progetto corrente
\item Recupera tutte le issue del progetto non archiviate
\item Per ogni issue recupera: ID, titolo, tipo, priorità, stato, autore, assegnatario, data
\item Applica eventuali filtri e ordinamenti salvati in sessione
\end{enumerate}

\paragraph{Output}
Lista issue con: tabella o card, icona tipo, titolo cliccabile, priorità (badge colorato), stato, assegnatario, pulsante "Crea Nuova Issue".

\end{reqbox}

\vspace{0.5cm}

\newpage

% =====================================================================
% REQ-FUN-11
% =====================================================================

\begin{reqbox}[title=REQ-FUN-11: Visualizza Dettagli Issue]
\label{subsubsec:req-visualizza-dettagli-issue}

\paragraph{Introduzione}
Gli utenti visualizzano il dettaglio completo di una issue del progetto corrente.

\paragraph{Precondizioni}
Utente autenticato, membro del progetto della issue.

\paragraph{Input}
Click su issue dalla lista.

\paragraph{Elaborazione}
\begin{enumerate}[leftmargin=1.1em]
\item Verifica che la issue appartenga al progetto corrente
\item Verifica che l'utente sia membro del progetto
\item Recupera tutti i dettagli della issue: informazioni base, commenti, etichette, immagini
\end{enumerate}

\paragraph{Output}
Pagina dettaglio / Modale con: header (ID, titolo, tipo, priorità, stato), informazioni (autore, data creazione, assegnatario), descrizione formattata, etichette, immagini (galleria), sezione commenti, pulsanti "Modifica" e "Archivia" (se amministratore).

\textbf{Errore}: "Issue non trovata" o "Accesso negato" se progetto non autorizzato.

\end{reqbox}

\vspace{0.5cm}

\newpage

% =====================================================================
% REQ-FUN-12
% =====================================================================

\begin{reqbox}[title=REQ-FUN-12: Modifica Issue]
\label{subsubsec:req-modifica-issue}

\paragraph{Introduzione}
Gli utenti possono modificare gli attributi delle issue del progetto corrente.

\paragraph{Precondizioni}
Utente autenticato, membro del progetto, issue non archiviata.

\paragraph{Input}
Modifiche a: titolo, descrizione, tipo, priorità, stato, etichette.

\paragraph{Elaborazione}
\begin{enumerate}[leftmargin=1.1em]
\item Verifica appartenenza al progetto
\item Valida tutti i campi modificati
\item Confronta i valori con quelli precedenti
\item Aggiorna la issue nel database
\item Se priorità diventa "Critica": notifica amministratori
\item Se stato diventa "Risolta": notifica autore della issue
\end{enumerate}

\paragraph{Output}
\textbf{Successo}: "Issue aggiornata con successo"

\end{reqbox}

\vspace{0.5cm}

\newpage

% =====================================================================
% REQ-FUN-13
% =====================================================================

\begin{reqbox}[title=REQ-FUN-13: Aggiungi Etichette]
\label{subsubsec:req-aggiungi-etichette}

\paragraph{Introduzione}
Gli utenti possono associare etichette personalizzate alle issue per facilitarne la categorizzazione.

\paragraph{Precondizioni}
Utente autenticato, membro del progetto.

\paragraph{Input}
Testo etichetta.

\paragraph{Elaborazione}
\begin{enumerate}[leftmargin=1.1em]
\item Normalizza il testo dell'etichetta
\item Verifica che non sia già associata alla issue
\item Verifica che il totale etichette non superi 10
\item Associa l'etichetta alla issue
\end{enumerate}

\paragraph{Output}
\textbf{Aggiunta}: Tag colorato visibile immediatamente.

\textbf{Rimozione}: Click su X del tag rimuove l'etichetta dalla issue.

\end{reqbox}

\vspace{0.5cm}

\newpage

% =====================================================================
% REQ-FUN-14
% =====================================================================

\begin{reqbox}[title=REQ-FUN-14: Aggiungi Immagine]
\label{subsubsec:req-aggiungi-immagine}

\paragraph{Introduzione}
Gli utenti possono allegare immagini alle issue per fornire documentazione visuale.

\paragraph{Input}
File immagine PNG o JPG.

\paragraph{Elaborazione}
\begin{enumerate}[leftmargin=1.1em]
\item Valida formato file e dimensione
\item Genera nome file univoco
\item Carica il file su storage
\item Associa l'immagine alla issue nel database
\end{enumerate}

\paragraph{Output}
\textbf{Successo: } "Immagine caricata con successo", pulsante "Rimuovi".

\textbf{Errore}: "Formato non supportato. Usa PNG o JPG"

\end{reqbox}

\vspace{0.5cm}

\newpage

% =====================================================================
% REQ-FUN-15
% =====================================================================

\begin{reqbox}[title=REQ-FUN-15: Filtra Issue]
\label{subsubsec:req-filtra-issue}

\paragraph{Introduzione}
Gli utenti possono applicare filtri per visualizzare sottoinsiemi specifici di issue del \textbf{progetto corrente}.

\paragraph{Input}
Selezione criteri di filtro: Tipo, Stato, Priorità, Etichetta, Assegnatario, Autore, Data creazione (range).

\paragraph{Elaborazione}
\begin{enumerate}[leftmargin=1.1em]
\item Costruisce i criteri di filtro selezionati dall'utente
\item Recupera dal database solo le issue del progetto corrente che soddisfano tutti i criteri
\item Salva i filtri applicati nella sessione per persistenza
\end{enumerate}

\paragraph{Output}
Lista issue aggiornata con solo gli elementi filtrati, badge "X filtri attivi".

\end{reqbox}

\vspace{0.5cm}

\newpage

% =====================================================================
% REQ-FUN-16
% =====================================================================

\begin{reqbox}[title=REQ-FUN-16: Ordina Issue]
\label{subsubsec:req-ordina-issue}

\paragraph{Introduzione}
Gli utenti possono ordinare la lista issue secondo criteri specificati.

\paragraph{Input}
Campo di ordinamento (Data creazione, Priorità, Stato, Titolo) e direzione (Crescente o Decrescente).

\paragraph{Elaborazione}
\begin{enumerate}[leftmargin=1.1em]
\item Applica l'ordinamento richiesto alla lista issue
\item Per Priorità usa ordine: Critica, Alta, Media, Bassa
\item Per Stato usa ordine: ToDo, In Lavorazione, In Revisione, Risolta, Chiusa
\item Salva la preferenza di ordinamento nella sessione
\end{enumerate}

\paragraph{Output}
Lista issue riordinata secondo il criterio selezionato, indicatore visuale (freccia) sulla colonna ordinata, persistenza durante la navigazione.

\end{reqbox}

\vspace{0.5cm}

\newpage

% =====================================================================
% REQ-FUN-17
% =====================================================================

\begin{reqbox}[title=REQ-FUN-17: Visualizza Issue Assegnate]
\label{subsubsec:req-visualizza-issue-assegnate}

\paragraph{Introduzione}
Gli utenti visualizzano una vista dedicata con solo le issue del \textbf{progetto corrente} assegnate a loro.

\paragraph{Input}
Click su "Le Mie Issue" dal menu.

\paragraph{Elaborazione}
\begin{enumerate}[leftmargin=1.1em]
\item Recupera tutte le issue del progetto corrente assegnate all'utente
\item Esclude le issue archiviate
\end{enumerate}

\paragraph{Output}
Lista filtrata con solo le issue assegnate all'utente, organizzazione per stato (ToDo, In Lavorazione, In Revisione, Risolta, Chiusa), possibilità di cambiare stato direttamente.

\end{reqbox}

\vspace{0.5cm}

\newpage

% =====================================================================
% REQ-FUN-18
% =====================================================================

\begin{reqbox}[title=REQ-FUN-18: Aggiungi Commento]
\label{subsubsec:req-aggiungi-commento}

\paragraph{Introduzione}
Gli utenti possono aggiungere commenti testuali alle issue per discussione e collaborazione.

\paragraph{Input}
Testo commento (obbligatorio).

\paragraph{Elaborazione}
\begin{enumerate}[leftmargin=1.1em]
\item Valida lunghezza testo
\item Sanitizza l'input
\item Crea nuovo commento associato alla issue con timestamp e autore
\item Se il commento menziona altri utenti (@username): invia notifiche
\item Aggiorna timestamp ultima modifica della issue
\end{enumerate}

\paragraph{Output}
Commento appare immediatamente nella sezione commenti con avatar autore, nome, timestamp, testo formattato, pulsante "Elimina" (visibile solo all'autore), toast "Commento aggiunto"

\end{reqbox}

\vspace{0.5cm}

\newpage

% =====================================================================
% REQ-FUN-19
% =====================================================================

\begin{reqbox}[title=REQ-FUN-19: Elimina Commento]
\label{subsubsec:req-elimina-commento}

\paragraph{Introduzione}
Gli utenti possono eliminare i propri commenti; gli amministratori possono eliminare qualsiasi commento.

\paragraph{Precondizioni}
Autore del commento oppure amministratore.

\paragraph{Elaborazione}
\begin{enumerate}[leftmargin=1.1em]
\item Mostra dialog conferma: "Eliminare commento? Operazione irreversibile."
\item Se confermato: elimina il commento dal database o lo marca come eliminato
\end{enumerate}

\paragraph{Output}
Commento scompare dalla lista o mostra "[Commento eliminato]", toast "Commento eliminato"

\end{reqbox}

\vspace{0.5cm}

\newpage

% =====================================================================
% REQ-FUN-19
% =====================================================================

\begin{reqbox}[title=REQ-FUN-19: Visualizza Notifiche]
\label{subsubsec:req-visualizza-notifiche}

\paragraph{Introduzione}
Gli utenti visualizzano notifiche relative ai progetti a cui hanno accesso.

\paragraph{Input}
Click sull'icona campanella nell'header.

\paragraph{Elaborazione}
\begin{enumerate}[leftmargin=1.1em]
\item Recupera tutte le notifiche dell'utente
\item Filtra per progetti a cui l'utente ha accesso
\item Ordina per timestamp decrescente (più recenti in alto)
\item Separa notifiche lette da non lette
\end{enumerate}

\paragraph{Output}
Dropdown o pagina notifiche con: lista notifiche (icona tipo, testo, timestamp relativo "2 ore fa"), badge rosso con numero notifiche non lette, evidenziazione notifiche non lette, pulsante "Segna tutte come lette"

\textbf{Tipi notifiche}: Issue assegnata, Issue Critica creata (solo admin), Issue risolta (per autore), Menzione in commento, Nuovo commento su issue seguita.

\end{reqbox}

\vspace{0.5cm}

\newpage

% =====================================================================
% REQ-FUN-20
% =====================================================================

\begin{reqbox}[title=REQ-FUN-20: Prendi Visione Notifica]
\label{subsubsec:req-prendi-visione-notifica}

\paragraph{Introduzione}
Gli utenti possono marcare le notifiche come lette.

\paragraph{Elaborazione}
\begin{enumerate}[leftmargin=1.1em]
\item Marca la notifica come letta nel database
\item Aggiunge timestamp di lettura
\item Decrementa il contatore badge notifiche non lette
\end{enumerate}

\paragraph{Output}
Notifica cambia aspetto (sfondo normale invece di evidenziato), badge contatore aggiornato, se cliccata reindirizza alla issue o commento relativo.

\end{reqbox}

\vspace{0.5cm}

\newpage

% =====================================================================
% REQ-FUN-21
% =====================================================================

\begin{reqbox}[title=REQ-FUN-21: Creazione Utenze (solo Admin)]
\label{subsubsec:req-creazione-utenze}

\paragraph{Introduzione}
Gli amministratori possono creare nuovi account utente.

\paragraph{Input}
Nome, Cognome, Email (univoca), Password (min 8 caratteri con requisiti di complessità), Ruolo (Utente o Amministratore).

\paragraph{Elaborazione}
\begin{enumerate}[leftmargin=1.1em]
\item Verifica ruolo amministratore
\item Valida tutti i campi e unicità email
\item Valida complessità password (almeno 1 maiuscola, 1 minuscola, 1 numero, 1 carattere speciale)
\item Crea nuovo account utente nel database
\end{enumerate}

\paragraph{Output}
\textbf{Successo}: "Utente [Nome Cognome] creato con successo", reindirizzamento a lista utenti

\textbf{Errore}: "Email già registrata", "Password non rispetta requisiti di complessità"

\end{reqbox}

\vspace{0.5cm}

\newpage

% =====================================================================
% REQ-FUN-22
% =====================================================================

\begin{reqbox}[title=REQ-FUN-22: Modifica Utenze (solo Admin)]
\label{subsubsec:req-modifica-utenze}

\paragraph{Introduzione}
Gli amministratori possono modificare i dati degli account esistenti.

\paragraph{Input}
Selezione utente, modifiche a: Nome, Cognome, Email, Password (opzionale), Ruolo.

\paragraph{Elaborazione}
\begin{enumerate}[leftmargin=1.1em]
\item Verifica che l'utente da modificare esista
\item Previene auto-modifica del ruolo amministratore corrente
\item Se email modificata: verifica unicità
\item Se password modificata: genera nuovo hash e invalida sessioni attive dell'utente
\item Aggiorna i dati dell'account nel database
\end{enumerate}

\paragraph{Output}
\textbf{Successo}: "Utente aggiornato con successo"

\textbf{Errore}: "Email già in uso", "Non puoi modificare il tuo ruolo amministratore"

\end{reqbox}

\vspace{0.5cm}

\newpage

% =====================================================================
% REQ-FUN-23
% =====================================================================

\begin{reqbox}[title=REQ-FUN-23: Rimozione Utenze (solo Admin)]
\label{subsubsec:req-rimozione-utenze}

\paragraph{Introduzione}
Gli amministratori possono eliminare account utente previa conferma.

\paragraph{Input}
Selezione utente da eliminare, password amministratore per conferma.

\paragraph{Elaborazione}
\begin{enumerate}[leftmargin=1.1em]
\item Dialog: "Eliminare utente? Operazione irreversibile. Le issue create rimarranno ma l'autore sarà marcato come 'eliminato'."
\item Verifica password amministratore
\item Previene auto-eliminazione
\item Invalida tutte le sessioni attive dell'utente
\item Marca le issue create dall'utente come "autore eliminato"
\item Rimuove l'utente da tutte le assegnazioni issue
\item Rimuove l'utente da tutti i progetti
\item Elimina l'account dal database
\end{enumerate}

\paragraph{Output}
\textbf{Successo}: "Utente [Nome Cognome] eliminato"

\textbf{Errore}: "Password amministratore non valida", "Non puoi eliminare il tuo account"

\end{reqbox}

\vspace{0.5cm}

\newpage

% =====================================================================
% REQ-FUN-24
% =====================================================================

\begin{reqbox}[title=REQ-FUN-24: Assegna Issue (solo Admin)]
\label{subsubsec:req-assegna-issue}

\paragraph{Introduzione}
Gli amministratori possono assegnare issue a membri del team del progetto.

\paragraph{Precondizioni}
Amministratore, issue e utente target esistenti, utente target membro del progetto della issue.

\paragraph{Input}
Issue da assegnare, selezione utente assegnatario dalla lista membri del progetto.

\paragraph{Elaborazione}
\begin{enumerate}[leftmargin=1.1em]
\item Verifica che l'utente target sia membro del progetto
\item Aggiorna l'assegnatario della issue
\item Se la issue era già assegnata: notifica il precedente assegnatario della riassegnazione
\item Crea notifica per il nuovo assegnatario
\end{enumerate}

\paragraph{Output}
"Issue \#[ID] assegnata a [Nome]", issue appare in "Le Mie Issue" dell'assegnatario, notifica inviata.

\textbf{Errore}: "Utente non è membro del progetto"

\end{reqbox}

\vspace{0.5cm}

\newpage

% =====================================================================
% REQ-FUN-25
% =====================================================================

\begin{reqbox}[title=REQ-FUN-25: Archivia Issue (solo Admin)]
\label{subsubsec:req-archivia-issue}

\paragraph{Introduzione}
Gli amministratori possono archiviare issue mantenendone la tracciabilità.

\paragraph{Precondizioni}
Amministratore, issue esistente non già archiviata.

\paragraph{Elaborazione}
\begin{enumerate}[leftmargin=1.1em]
\item Dialog: "Archiviare issue \#[ID]? Non sarà più visibile nelle liste ma rimarrà accessibile dall'archivio."
\item Se confermato: marca la issue come archiviata con timestamp e amministratore che ha eseguito l'operazione
\item La issue viene esclusa dalle liste standard
\item Rimane accessibile dalla vista "Archivio" (solo amministratori)
\end{enumerate}

\paragraph{Output}
"Issue \#[ID] archiviata", scompare dalla lista standard, pulsante "Ripristina" disponibile.

\textbf{Ripristino}: L'amministratore può ripristinare la issue dalla vista Archivio.

\end{reqbox}

\vspace{0.5cm}

\newpage

% =====================================================================
% REQ-FUN-26
% =====================================================================

\begin{reqbox}[title=REQ-FUN-26: Visualizza Report Mensile (solo Admin)]
\label{subsubsec:req-visualizza-report-mensile}

\paragraph{Introduzione}
Gli amministratori visualizzano report mensile del \textbf{progetto selezionato} con metriche aggregate.

\paragraph{Precondizioni}
Amministratore, progetto selezionato, dati di almeno un mese di attività.

\paragraph{Input}
Selezione mese e anno tramite date picker (default: mese corrente), click su "Genera Report".

\paragraph{Elaborazione}
Il sistema calcola le seguenti metriche per il progetto corrente nel periodo selezionato:

\textbf{Metriche Globali del Progetto}:
\begin{itemize}[leftmargin=1.1em]
\item Numero issue create e risolte nel mese
\item Numero issue ancora aperte
\item Distribuzione percentuale per tipo e priorità
\item Tempo medio di risoluzione in ore
\end{itemize}

\textbf{Metriche per Utente (membri del progetto)}:
\begin{itemize}[leftmargin=1.1em]
\item Per ogni membro: issue create, assegnate, risolte
\item Tempo medio risoluzione personale
\item Percentuale efficienza (risolte su assegnate)
\end{itemize}

\textbf{Allerte}: lista issue critiche non risolte da più di 7 giorni

\paragraph{Output}
Dashboard report con: card KPI con numeri principali, grafici per distribuzioni, tabelle ripelogative per utente.

\textbf{Errore}: "Permessi insufficienti", "Nessun dato disponibile per il periodo selezionato"

\end{reqbox}

\newpage
